\section{Actual analytic realization of normal weighted components}

\begin{definition}[The genuine normal homogeneous polynomial of an original analytic germ]
\label{hir:actual-normal-homogeneous-polynomial}
\uses{hir:actual-relative-coefficient-polynomial-truncations,hir:actual-relative-coefficient-germ-stability}
\lean{Hironaka.HironakaActualNormalHomogeneousComponents.actualNormalHomogeneousPolynomial,Hironaka.HironakaActualNormalHomogeneousComponents.actualNormalHomogeneousPolynomial_coefficient,Hironaka.HironakaActualNormalHomogeneousComponents.actualNormalHomogeneousPolynomial_isHomogeneous}
\leanok
For an original ambient analytic germ $g$ in $d+t$ coordinates and
$q\ge0$, take the degree-$q$ homogeneous component of its genuine
normal Taylor polynomial modulo the $(q+1)$-st normal power. This is a
finite polynomial in the normal variables with original entire tangent
holomorphic germs as coefficients. Its coefficient at $\gamma$ is the
actual stable relative coefficient of $g$ if $|\gamma|=q$, and zero
otherwise. It is homogeneous of normal degree $q$.
\end{definition}
\begin{proof}
The normal Taylor polynomial is constructed by the inverse of the
actual normal-power quotient basis. Its coefficients of degree below
$q+1$ are the genuine stable relative coefficient germs. The original
finite polynomial homogeneous-component map retains exactly coefficients
of degree $q$, all of which lie in that range, and replaces all other
coefficients by zero. Its genuine homogeneous-component theorem gives
the stated normal homogeneity. No infinite series is evaluated away
from its origin.
\end{proof}

\begin{theorem}[The actual analytic normal symbol realizes the entire formal weighted component]
\label{hir:actual-normal-homogeneous-entire-taylor-realization}
\uses{hir:actual-normal-homogeneous-polynomial,hir:actual-relative-polynomial-formal-taylor-coefficient,hir:actual-relative-coefficient-literal-formal-taylor-dictionary,hir:actual-weighted-component-definition}
\lean{Hironaka.HironakaActualNormalHomogeneousComponents.actualTaylor_normalHomogeneousPolynomial}
\leanok
Evaluate that genuine finite normal polynomial using the original
normal coordinate germs and actual tangent pullback. Its entire original
Taylor series is exactly the actual weight-$q$ component of the entire
Taylor series of $g$, with weights zero on the tangent block and one
on the normal block.
\end{theorem}
\begin{proof}
Every actual exponent splits uniquely into its original tangent exponent
$\beta$ and normal exponent $\gamma$. The genuine relative polynomial
Taylor dictionary identifies its joined coefficient with the $\beta$
coefficient of the true tangent Taylor series of the polynomial's
$\gamma$ coefficient. If $|\gamma|=q$, the preceding coefficient formula
and the literal whole relative-coefficient dictionary identify this
with the original joined coefficient of Taylor$(g)$. Otherwise it is
zero. The actual weight of a joined exponent is precisely
$|\gamma|$, so this is exactly the coefficient rule for the original
weighted component at every exponent. Genuine formal coefficient
extensionality proves equality of the entire series.
\end{proof}

\begin{theorem}[The entire analytic vertex-symbol ideal has precisely the genuine weighted formal extension]
\label{hir:actual-analytic-normal-vertex-symbol-formal-extension}
\uses{hir:actual-normal-homogeneous-entire-taylor-realization,hir:actual-analytic-vertex-coefficient-data}
\lean{Hironaka.HironakaActualNormalHomogeneousComponents.actualAnalyticNormalVertexComponentIdeal,Hironaka.HironakaActualNormalHomogeneousComponents.actualAnalyticNormalVertexComponentIdeal_formal_extension}
\leanok
For genuine analytic monic vertex data of the original ideal $I$, span
the actual normal homogeneous polynomial evaluations of the actual
vertex members, each at the normal weight of its true vertex exponent.
The entire original Taylor extension of this analytic ideal is exactly
the span of the genuine formal weighted components of those same
original Taylor members at those same vertex weights.
\end{theorem}
\begin{proof}
The original Taylor algebra homomorphism sends each actual analytic
generator to its displayed entire formal component by the preceding
theorem. The image of a whole ideal span is the whole span of the actual
images. Applying this literal identity to the actual generator range
gives both entire ideal directions, retaining the original coefficients,
vertex weights and monomial order.
\end{proof}

\begin{theorem}[Only the original analytic ideal supplies every whole analytic vertex-symbol HS bound]
\label{hir:actual-original-analytic-normal-vertex-symbol-hs-bound}
\uses{hir:actual-original-analytic-vertex-coefficient-data-existence,hir:actual-analytic-normal-vertex-symbol-formal-extension,hir:actual-weighted-vertex-ideal-hs-bound,hir:actual-unshifted-original-formal-hs-transfer}
\lean{Hironaka.HironakaActualNormalHomogeneousComponents.exists_actualAnalyticNormalVertexComponentIdeal_hs_le}
\leanok
Every original analytic ideal internally supplies genuine analytic monic
vertex data for which the entire actual normal-symbol ideal has every
unchanged original analytic Hilbert--Samuel value at most the
corresponding unchanged original value of $I$.
\end{theorem}
\begin{proof}
Construct genuine analytic vertex coefficient data from the original
ideal using the proved existence theorem. Their actual Taylor members
belong to the genuine whole formal extension, have coefficient one at
the actual vertices and vanish below them, hence give a genuine monic
vertex family for that entire original formal ideal. The preceding
whole ideal image identity identifies the actual analytic symbol ideal's
formal extension with its genuine weighted vertex-component ideal.
The proved weighted diagram inclusion and standard-index injection give
the all-order genuine formal HS inequality. The actual whole Taylor
quotient equivalence identifies each of these genuine formal dimensions
with the unchanged original analytic HS value. Combining these facts
proves every order simultaneously, with no coefficient, convergence,
comparison or finite-generator certificate supplied by the caller.
\end{proof}

This constructs actual analytic normal symbols and their full HS bound.
Their comparison with actual exceptional normalized quotients still
requires genuine centre-power compatibility and the true exceptional
restriction calculation. Selection of general permissible centres and
the general positive-dimensional Bennett theorem remain open.
